\scanpage{14}
\deleted{\unsure{Comme}}
\deleted{Soit donc $Y$ une hypersurface \deleted{\unsure{non irréductible}}\inserted{à composantes disjointes} dans $X$.
\inserted{$x^1\in A^1(X)$ sa classe}
On a
\[
\widetilde C_X(Y)=(1-\deleted{\unsure{d(Y)}}x^1)^{-1}=\sum x^{1n}
\]
\[
\begin{aligned}
\widetilde C_{X\times P^r}(i_!(\underline O_Y))
&=\widetilde C_{X\times P^r}(Y\times P^r)\\
&=\deleted{1+(-1)^{r-1}(r-1)!\sum_{i=1}^\infty Q_i^{(r)}(x^1,(x^1)^2\text{---}\unsure{x^i})}\\
&\quad1+(-1)^{r-1}(r-1)!\eta^r\sum_{i=1}^\infty
\underbrace{Q_i^{(r)}(x^1,(x^1)^2,\text{---}\,(x^1)^{\unsure{r|i}})}_{C_{i,r}(x^1)^{\unsure{r|i}}}
\end{aligned}
\]
\deleted{$1+\sum$}
Si donc $(r-1)!X^1=0$, on a $\widetilde C_{X\times P^r}(Y\times P^r)=0$,
sans que pourtant l'on ait nécessairement $X^1=0$.
Cas le plus simple : $r=3$, $X$ une courbe elliptique.}

\uline{N.B.} Il se peut qu'on ait un élément $x$ de $\underline K(X)$ tel que les
$\unsure{C^i|c^i}(x)$ sont tous annulés par $(r-1)!$, sans que $x$ soit nul.
Alors $y=i_!(x)$ a la propriété que $\widetilde C(\deleted{\unsure{x}}\inserted{y})=0$,
mais $y\ne0$. Exemple le plus simple :

$X=$ courbe elliptique, $\inserted{x\in\underline K(X)}$ \unsure{classe} de diviseur \unsure{non}
nul d'ordre $2$, [p.ex. $x=(a)-(0)$, où $a\in X$ est tel que $a+a=0$, $a\ne0$], et $r=3$.
\uline{On a donc dans $P^3\times X$ la différence de deux $\unsure{V^0|V^s}$
\deleted{\unsure{points}} dont la classe de Chern est nulle ; sans qu'elle soit nulle.}

\margin{13}
